\chapter{Music Datasets}\label{chapter:music_datasets}
  Since this project deals with music genre classification, a suitable set of musical stimuli have to be chosen to gather valuable data. As research in the domain of music is widespread, there exist multiple common music datasets used in different studies. Therefore, the difficulty lies in choosing the right dataset for this thesis. However, the project goal cristallizes out clear requirements for such a music dataset.\par
  First of all, the songs that will be listened to are suppossed to be around 30s long. The reason for that is that the stimulus could go lost if the song runs too long because the brain can adapt to the genre-specific features, thus reducing the visibility of the impact on the brain. Apart from that, the amount of songs should be quite large with around 800 to 1000 song clips. It is important to have a large pool of songs to get as many different and versatile data as possible. This helps in future steps regarding classification. However, related to this requirement, the genres used in this pool also have to meet specifications. First of all, the songs must be clearly labeled to avoid confusion in later stages of the study. This is also important to know for sure how balanced the dataset is in terms of genres listened to. This is supposed to be as balanced as possible so that each genre is represented equally. Additionally, it should be easy to employ this dataset in the study. It should be uncomplicated to load the songs as well as getting access to labels or other important information.\par
  In the following, different popular and public music datasets are presented. They are going to be analyzed according their possible use in this research and on how they meet the requirements specified above. After that, the reasons for the final choice of dataset are outlined.
  \section{Audio Set}
    In 2017, Gemmeke et al. published a dataset for audio events called Audio Set\footnote{\url{https://research.google.com/audioset/}}. This is a large-scale dataset consisting of audio events labeld by humans using a carefully structured hierarchical ontology of 632 audio classes. The source of the audio events are derived from 10 second segments of YouTube videos. These audio events are occurences in the real-world, like human sounds, animal sounds, natural sounds, sounds of things, source-ambigious sounds, environment sounds and music. Then, each of these sounds is labeled with one ontology and labels with a maximal hierarchical depth of six. In the determination of labels, the researchers made sure that the tags are intuitive and clearly distinct from each other. \cite{gemmeke2017}\par
    In terms of the use for this project, the audio set has some drawbacks. First of all, the focus of this dataset does not lie in music but in general audio events. This is why in the first step, the dataset would have to be filtered to just contain the music clips. However, the fact that the clips are only 10 seconds long, there is the risk that this is not enough of a stimulus to arouse the brain and see the distinct patterns. Another problem lies in the format of the data. Since every audio event is a YouTube link, each sound would have to be downloaded from that link to allow a seamless integration with the experiment setup, thus illustrating a large overhead. Last but not least, the dataset is not labeled in regards of specific genres.\textcolor{red}{WRONG CHANGE THIS: SEE THIS PAPER FOR REFERENCE: \url{https://arxiv.org/pdf/1804.01149} ; THEY HAVE USED IT FOR GENRE CLASSIFICATION } Therefore this not only adds the overhead of manually mach each music piece to its genre but also does not guarantee that each genres appears the same amount of times in the dataset, making it a not a so well fit for the grounds of this thesis.

  \section{Million Song Dataset}
    The Million Song Dataset\footnote{\url{http://millionsongdataset.com/}} is a collection of audio features and metadata for a million music tracks, which also is freely available. At the moment, this is the largest research dataset in the field as it contains 1,000,000 songs accumulating around 280GB of data. In addition to the songs mainly from western commercial music in 2011, general metadata and audio features like pitch, timbre, etc. are provided by making use of the EchoNest API. This API is an online resource supporting many music applications on the web. Since most of the "natively" available metadata are coming from the EchoNest API, it is easily extendable to further labels like genres. \cite{bertin2011million}\par
    Although this dataset may seem like a great fit because of its rich amount of songs, the amount of songs is too much. Since only 800-1000 songs are used during the acquisition, a majority of the dataset would have to be filtered out. Furthermore, all songs are not trimmed down to a shorter clip but are provided in their full length, which is not meeting the requirement for this study. Furthermore, the genres would have to be added externally since they are not available in the native dataset. This also emphasizes that it does not guarantee any balance in genres available over the course of the 1,000,000 songs.

  \section{MagnaTagATune}
    One challenge that comes up when creating such datasets is how a label or tag is determined for each song. In the context of genres, how is it decided which genre which song belongs to, as there could be overlapping features. On the basis of this hurdle, Law et al. created a music dataset to evaluate different tagging algorithms called MagnaTagATune\footnote{\url{https://paperswithcode.com/dataset/magnatagatune}}. It contains 25,863 music clips each 29 seconds long. These excerpts were taken from 5223 songs. These music pieces range across many different genres like Rock, Pop, Jazz, Blues, World, Metal, Punk and many more. In general, each clip is accompanied by a vector of binary annotations of 188 tags using the TagATune game. \cite{güçlü2016brainsbeats}\par
    TagATune is a game with two players that aims to collect music tags from the players. Each player is presented with a music clip which they are assigned to give them labels. After that, both players see what the other player has chosen as tags. Based on this information, the players have to decide whether they listened to the same music clip. \cite{law2009evaluation}\par
    In perspective of this study, this music dataset has the advantage that each clip is around 30 seconds long, hitting the desired length for one stimulus. Furthermore, it is beneficial that the songs vary in genres used, thus being a good candidate as a music dataset. However, there are some drawbacks, especially lying in the labeling of each clip. Although the labeling process is well-thought and transparent, each song can have more than one label. This means that, depending on the players, one clip can be annotated with more than one genre tag. And since there are in total 37 genre related tags \cite{rotheneder2023explanation}, deciding which genre is the most suitable choice for a song is very difficult. Related to that, since there are around 26,000 clips, one would have to remove majority of songs in order to let it fit into the scope of this thesis. Combined with the fact that multiple genres can be chosen for each song, it is difficult to achieve a balanced dataset according in regards of genre.

  \section{FMA Dataset}
    The \ac{FMA}\footnote{\url{https://github.com/mdeff/fma}} is an open dataset designed for multiple tasks in the domain of \ac{MIR}
  \section{GTZAN Genre classification}

  \section{Final decision}
    \textcolor{orange}{Add a table summarizing the datasets. As columns, it should have name, amount of songs/music clips, average length of songs, amount of genres, availability of audio files}